% Chapter 13: The Virtual Haken Conjecture and the Surface Subgroup Conjecture
\let\LocalMacro\relax

\chapter{The Virtual Haken Conjecture and the Surface Subgroup Conjecture}

\section{The Problem}

\subsection{Informal}
Imagine you have a tangled, complicated three-dimensional shape---like a knotted ball of yarn frozen in space. Can you always find a surface tucked inside that helps you understand the shape's structure? Think of it as a soap film that cuts cleanly through the interior, revealing how everything fits together. For decades, mathematicians wondered whether every complicated 3D shape must contain such a helpful surface.

\subsection{Formal}
\begin{conjecture}[Virtual Haken Conjecture, Waldhausen \cite{thurston1982}]
Every compact, orientable, irreducible 3-manifold with infinite fundamental group has a finite cover that is Haken.
\end{conjecture}

A stronger version is the \emph{Surface Subgroup Conjecture}:

\begin{conjecture}[Surface Subgroup Conjecture]
If $M$ is a closed hyperbolic 3-manifold, then $\pi_1(M)$ contains a subgroup isomorphic to the fundamental group of a closed surface of genus $g \ge 2$.
\end{conjecture}

\begin{historicalbox}
Waldhausen asked these questions in the 1960s. They became central to 3-manifold topology: the Virtual Haken Conjecture would reduce the study of general 3-manifolds to Haken manifolds (which are well-understood via normal surface theory). Thurston's geometrization program made them central to the classification of 3-manifolds.
\end{historicalbox}

\section{Mathematical Theory}


\subsection{Prerequisites}
This chapter assumes familiarity with:
\begin{itemize}
\item 3-manifold topology (fundamental group, covering spaces)
\item Hyperbolic geometry (Thurston's geometrization)
\item Geometric group theory (quasiconvexity, CAT(0) spaces)
\item Basic low-dimensional topology (incompressible surfaces, Haken manifolds)
\end{itemize}

\subsection{Haken Manifolds and Normal Surfaces}
A manifold is \emph{irreducible} if every embedded 2-sphere bounds a 3-ball. A surface $S \subset M$ is \emph{incompressible} if the map $\pi_1(S) \to \pi_1(M)$ is injective.

Haken's algorithm for the homeomorphism problem uses \emph{normal surfaces}: surfaces that intersect each tetrahedron of a triangulation in a standard way. The space of normal surfaces forms a cone in $\mathbb{R}^{7t}$ (where $t$ is the number of tetrahedra), and fundamental solutions correspond to incompressible surfaces.

\subsection{The Virtually Fibered Conjecture}
Thurston conjectured an even stronger statement:

\begin{conjecture}[Virtually Fibered Conjecture]
Every closed hyperbolic 3-manifold has a finite cover that fibers over the circle.
\end{conjecture}

If a manifold fibers over the circle, it is Haken (the fiber is an incompressible surface). So Virtual Fibering $\implies$ Virtual Haken $\implies$ Surface Subgroup.

\section{The Solution}

\subsection{Kahn and Markovic \cite{kahnmarkovic2012}: Surface Subgroups}
Jeremy Kahn and Vladimir Markovic (2012) proved the Surface Subgroup Conjecture using:

\begin{theorem}[Kahn-Markovic \cite{kahnmarkovic2012}]
Every closed hyperbolic 3-manifold contains an immersed, quasi-Fuchsian surface subgroup.
\end{theorem}

Their method:
\begin{enumerate}
\item \textbf{Good pants}: Build pairs of pants with cuff lengths $\approx L$ and twists $\approx 0$ in the hyperbolic manifold.
\item \textbf{Exponential mixing}: Use the exponential mixing of the geodesic flow to glue pants together with nearly matching cuff lengths.
\item \textbf{Almost isometry}: The glued surface is $\epsilon$-close to a totally geodesic surface, giving a quasi-Fuchsian embedding.
\end{enumerate}

\begin{highlightbox}
Kahn and Markovic's proof was a masterpiece of ergodic theory applied to 3-manifolds. They used the exponential mixing of the geodesic flow (which they proved) to find almost-perfectly-matching pants. The proof is constructive but non-algorithmic.
\end{highlightbox}

\subsection{Agol \cite{agol2013}: Virtual Haken and Virtual Fibering}
Ian Agol \cite{agol2013} proved the full Virtual Haken Conjecture and Virtual Fibering Conjecture using:

\begin{theorem}[Agol \cite{agol2013}]
Every closed hyperbolic 3-manifold is virtually fibered (and hence virtually Haken).
\end{theorem}

Agol's proof combined:
\begin{enumerate}
\item \textbf{Wise's work on cubulation}: Daniel Wise \cite{wise2012} showed that the fundamental group of a hyperbolic 3-manifold acts geometrically on a CAT(0) cube complex.
\item \textbf{Agol's criterion}: Agol proved that if a hyperbolic group acts geometrically on a CAT(0) cube complex, then it is virtually special (has a finite-index subgroup embedding into a right-angled Artin group).
\item \textbf{Virtual fibering}: Virtual specialness implies virtual fibering (by work of Agol and Wise).
\end{enumerate}

\begin{theorem}[Virtual Haken Conjecture, Agol \cite{agol2013}]
Every compact, orientable, irreducible 3-manifold with infinite fundamental group has a Haken finite cover.
\end{theorem}

\begin{highlightbox}
Agol's proof was a synthesis of 3-manifold topology, geometric group theory, and CAT(0) cube complex theory. The key insight was that hyperbolic 3-manifold groups are "cubulated" by Wise, and Agol showed this forces virtual fibering.
\end{highlightbox}

\section{Impact}
\begin{itemize}
\item \textbf{3-manifold classification}: Reduced the study of general 3-manifolds to Haken manifolds (solvable by normal surface theory).
\item \textbf{Geometric group theory}: Proved that hyperbolic 3-manifold groups are linear, virtually special, and have many surface subgroups.
\item \textbf{Cubical geometry}: Established CAT(0) cube complexes as a central tool in 3-manifold topology.
\item \textbf{Quantum topology}: Implications for the volume conjecture and TQFT.
\end{itemize}

\section{Exercises}

\begin{exercise}[Basic]
Define incompressible surface. Why is a sphere incompressible in $S^3$ but compressible in $S^2 \times S^1$?
\end{exercise}

\begin{exercise}[Intermediate]
Explain why a manifold that fibers over the circle is Haken.
\end{exercise}

\begin{exercise}[Advanced]
Sketch how the exponential mixing of the geodesic flow is used to glue good pants into a quasi-Fuchsian surface.
\end{exercise}

\section{Timeline}

\begin{tabularx}{\linewidth}{lX}
\toprule
\textbf{Year} & \textbf{Event} \\ \midrule
1968 & Waldhausen poses Virtual Haken and Virtual Fibered conjectures \\ 
1982 & Thurston's geometrization program includes these as central questions \\ 
2002 & Perelman proves geometrization \\ 
2009 & Wise begins cubulation program for hyperbolic 3-manifolds \\ 
2012 & Kahn-Markovic prove Surface Subgroup Conjecture \\ 
2012 & Agol proves Virtual Haken and Virtual Fibering conjectures \\ 
2013 & Agol wins Breakthrough Prize; Kahn-Markovic win Clay Research Award \\ 
2016 & Agol wins Fields Medal \\ \bottomrule
\end{tabularx}

\section{Citations}

\begin{enumerate}
\item Agol, I. (2013). ``The virtual Haken conjecture.'' Documenta Mathematica, 18, 1045--1087.
\item Kahn, J., \& Markovic, V. (2012). ``Immersing almost geodesic surfaces in a closed hyperbolic three manifold.'' Annals of Mathematics, 175(3), 1127--1190.
\item Wise, D. T. (2012). ``The structure of groups with a quasiconvex hierarchy.'' Annals of Mathematics Studies, 209.
\item Agol, I., Groves, D., \& Manning, J. (2012). ``An alternate proof of Wise's malnormal special quotient theorem.'' Forum of Mathematics, Pi, 1, e1.
\item Thurston, W. (1982). ``Three-dimensional geometry and topology.'' Princeton University Press.
\end{enumerate}
